%+++++++++++++++++++++++++++++++++++++++++++++++++++++++++++++++++++++++++++++++++++++++++++++++++++++++++++++++++++++++++++
\section{Les complexes}
%+++++++++++++++++++++++++++++++++++++++++++++++++++++++++++++++++++++++++++++++++++++++++++++++++++++++++++++++++++++++++++

\begin{probleme}
    Encore une fois, cette section n'est pas du tout faite. Le but de cette section serait de
    \begin{itemize}
        \item Construire \( \eC\) en tant qu'ensemble
        \item Construire les opérations courantes.
        \item Démontrer les propriétés de base.
    \end{itemize}
    Attention : il est sans espoir de parler de forme trigonométrique ici parce que les exponentielles et fonctions trigonométriques ne sont définies qu'avec les séries.

    Nous donnons ici en vrac quelques propriétés et définitions que se doivent d'être présentes dans la version finale de cette section, si elle existe un jour.

    Comme partout dans ce chapitre sur la construction des ensembles de nombres, certains propriétés qui ont l'air toute simples peuvent, en fonction des définitions prises, s'avérer pas du tout évidentes.
\end{probleme}


 \subsection{Définitions}
 Un nombre complexe s'écrit sous la forme $z = a + b i$, où $a$ et $b$
 sont des nombres réels appelés (et notés) respectivement partie réelle
 ($a = \Re(z)$) et partie imaginaire ($b = \Im(z)$) de $z$. L'ensemble
 des nombres de cette forme s'appelle l'ensemble des nombres complexes
 ; cet ensemble porte une structure de corps et est noté $\eC$. Le
 nombre complexe $i = 0 + 1 i$ est un nombre imaginaire qui a la
 particularité que $i^2 = -1$.

 Deux nombres complexes $a + bi$ et $c + di$ sont égaux si et seulement
 si $a = c$ et $b = d$, c'est-à-dire leurs parties réelles sont égales,
 et leurs parties imaginaires sont égales.

 Pour $z = a + bi$ un nombre complexe, on note $\bar z = a - bi$ le
 \Defn{complexe conjugué} de $z$. Dans le plan de Gauss, il s'agit du
 symétrique de $z$ par rapport à la droite réelle (généralement
 dessinée horizontalement).

 On définit le module du complexe $z$ par $\module z = \sqrt{z\bar z} =
 \sqrt{a^2 + b^2}$. Dans le plan de Gauss, il s'agit de la distance
 entre $0$ et $z$.

\begin{proposition}
Pour tous $z = a+bi$ et $z^\prime$ nombres complexes, on a
   \begin{enumerate}
   \item $z \bar z = a^2 + b^2$;
   \item $\bar{\bar{z}} = z$;
   \item $\module z = \module {\bar z}$;
   \item $\module{zz^\prime} = \module z \module{z^\prime}$;
   \item $\module{z+z^\prime} \leq \module z + \module{z^\prime}$.
   \end{enumerate}
\end{proposition}

\begin{lemma}   \label{LEMooONLNooXLNbtB}
    Pour tout \( z\in \eC\) nous avons \( z\bar z=\bar z z=| z |^2\).
\end{lemma}
